\documentclass[11pt,a4paper]{article}
% =====================================================================
% preambulo.tex: modelo fixo dos relatórios de AEDS2 (PUCRS)
%
% Para um novo trabalho: copie a pasta inteira, troque título e autor em
% main.tex e reescreva as seções em secoes/. Este arquivo não precisa mudar.
%
% Compila com pdfLaTeX (padrão do Overleaf) e com XeLaTeX. O ramo \ifPDFTeX
% escolhe as fontes de cada motor; o resto é igual.
% =====================================================================

\usepackage{iftex}
\ifPDFTeX
  \usepackage[utf8]{inputenc}
  \usepackage[T1]{fontenc}
\else
  \usepackage{fontspec}
\fi
\usepackage[brazilian]{babel}
\usepackage{amsmath}
\ifPDFTeX
  \usepackage{newtxtext,newtxmath}   % Times no texto e na matemática
\else
  % Fontes chamadas pelo nome do arquivo: funciona no XeLaTeX completo e no Tectonic.
  \setmainfont{texgyretermes}[Extension=.otf,UprightFont=*-regular,BoldFont=*-bold,
    ItalicFont=*-italic,BoldItalicFont=*-bolditalic]
  \setmonofont{texgyrecursor}[Extension=.otf,UprightFont=*-regular,BoldFont=*-bold,
    ItalicFont=*-italic,BoldItalicFont=*-bolditalic,Scale=MatchLowercase]
\fi

% ---------------------------------------------------------------- página e texto
\usepackage[a4paper,margin=2.5cm]{geometry}
\setlength{\parindent}{0pt}
\setlength{\parskip}{0.6em}
\emergencystretch=1.5em

\usepackage{enumitem}
\setlist{itemsep=0.25em,topsep=0.3em}

% ---------------------------------------------------------------- cores, figuras, gráficos
\usepackage{xcolor}
\usepackage{graphicx}
\usepackage{tikz}
\usetikzlibrary{arrows.meta}
\usepackage{pgfplots}
\pgfplotsset{compat=1.17}

% ---------------------------------------------------------------- tabelas e legendas
\usepackage{array}
\renewcommand{\arraystretch}{1.2}
\usepackage[font=small,labelfont=it,textfont=it,labelsep=endash,justification=centering]{caption}
\captionsetup{skip=6pt}

% ---------------------------------------------------------------- pseudo-código numerado
% Linhas numeradas permitem citar "linha 7" no texto e ligar cada etapa à explicação.
\usepackage{algorithm}
\usepackage{algpseudocode}
\floatname{algorithm}{Algoritmo}
\algrenewcommand\algorithmicreturn{\textbf{devolve}}
\algrenewtext{Function}[2]{\textbf{função} \textproc{#1}(#2)}
\algrenewtext{EndFunction}{\textbf{fim função}}
\algrenewtext{While}[1]{\textbf{enquanto} #1 \textbf{faça}}
\algrenewtext{EndWhile}{\textbf{fim enquanto}}
\algrenewtext{If}[1]{\textbf{se} #1 \textbf{então}}
\algrenewtext{EndIf}{\textbf{fim se}}
\algrenewtext{For}[1]{\textbf{para} #1 \textbf{faça}}
\algrenewtext{ForAll}[1]{\textbf{para cada} #1 \textbf{faça}}
\algrenewtext{EndFor}{\textbf{fim para}}

% ---------------------------------------------------------------- PDF
\usepackage[hidelinks]{hyperref}

% ---------------------------------------------------------------- macros do modelo
% \pendente{texto}: lembrete em vermelho para resolver antes de entregar.
% Antes da entrega, procure por "pendente" nos arquivos .tex e apague todos.
\newcommand{\pendente}[1]{\textcolor{red}{\textbf{[#1]}}}

% \espacofigura{altura}{indicação}: caixa tracejada que reserva o lugar de uma figura
% e diz o que desenhar. Quando a imagem existir, troque a chamada por
% \includegraphics[width=0.9\linewidth]{figuras/arquivo}.
\newcommand{\espacofigura}[2]{%
  \begin{tikzpicture}
    \node[draw=gray,dashed,rounded corners=2pt,text width=0.86\linewidth,
          minimum height=#1,align=center,inner sep=6pt,font=\small\itshape]
      {\textbf{Espaço para a figura}\\[3pt] #2};
  \end{tikzpicture}}


\hypersetup{pdftitle={Trabalho 2: O Cavalo e os Hiperpulos},
            pdfauthor={Guilherme Guimarães Santos}}

% Relatórios não têm capa: título, autor e data ficam no topo da primeira página.
\title{Trabalho 2: O Cavalo e os Hiperpulos}
\author{Guilherme Guimarães Santos\\[2pt]
  \small Escola Politécnica, PUCRS\\
  \small Algoritmos e Estruturas de Dados II, Turma 320\\
  \small Prof. João Batista Souza de Oliveira}
\date{11 de novembro de 2026}

\begin{document}
\maketitle

\begin{abstract}
Este trabalho trata de um desafio de xadrez com regras diferentes das habituais: em um
tabuleiro quadrado cujas bordas se encostam, cada casa traz um dígito de 0 a 9 que aumenta
o pulo em L do cavalo, e o cavalo deve ir da casa C até a saída S no menor número de
movimentos. O tabuleiro é modelado como um grafo, com uma casa por vértice e um pulo por
aresta, e o problema é resolvido com busca em largura (BFS), que encontra o menor caminho
quando todas as arestas custam o mesmo. O custo da BFS é proporcional ao número de casas,
$\Theta(N^2)$ para um tabuleiro $N \times N$, e as medições confirmam: dobrar $N$ multiplica
o tempo por cerca de quatro. Os oito tabuleiros fornecidos, de $40 \times 40$ a
$1500 \times 1500$, foram resolvidos com 4 a 43 movimentos, cada um em no máximo 57 ms, e as
respostas foram conferidas por uma segunda implementação, por busca exaustiva e pelo
algoritmo de Dijkstra.
\end{abstract}

\section{Introdução}
\label{sec:introducao}

O problema tratado neste trabalho é um desafio de xadrez com regras diferentes das habituais.
Dado um tabuleiro quadrado de lado $N$, com um dígito de 0 a 9 em cada casa, e um cavalo em
uma casa marcada com C, o objetivo é levar o cavalo até uma casa de saída, marcada com S, no
menor número possível de movimentos.

Duas regras tornam o problema diferente do xadrez comum. A primeira é que o dígito das casas
aumenta o pulo do cavalo: o dígito 0 corresponde ao pulo normal de xadrez, em forma de L, e
dígitos maiores esticam o L, como mostra a Figura~\ref{fig:pulos}. A segunda é que as bordas
do tabuleiro se encostam: se o cavalo sai pela direita, entra pela esquerda, e se sai por
cima, entra por baixo. Por isso diz-se que o tabuleiro é \emph{toroidal}.

\begin{figure}[ht]
  \centering
  % Figura 1: forma do pulo em L para os dígitos 0, 1 e 2.
%
% Como ler o código: cada L tem (2+d) passos na horizontal e (1+d) na vertical.
% As casas brancas formam o L do dígito 0; as casas coloridas são as d casas
% acrescentadas em cada perna. Para mudar as cores, edite \colorlet abaixo; para mostrar
% mais dígitos, acrescente um par "d/posição" na lista do \foreach (por exemplo 3/17.2).
\begin{tikzpicture}[x=0.62cm,y=0.62cm,font=\footnotesize]
  \colorlet{extra}{orange!55}
  \foreach \d/\xo in {0/0, 1/4.6, 2/10.2} {
    \pgfmathtruncatemacro{\pl}{2+\d}   % passos na horizontal = coluna do canto
    \pgfmathtruncatemacro{\pc}{1+\d}   % passos na vertical
    % perna horizontal (linha de cima)
    \foreach \c in {0,...,\pl} {
      \pgfmathtruncatemacro{\ex}{and(\c>=2,\c<\pl)}
      \ifnum\ex=1 \def\cor{extra}\else\def\cor{white}\fi
      \draw[fill=\cor] ({\xo+\c},0) rectangle ++(1,1);
    }
    % perna vertical (abaixo do canto)
    \foreach \r in {1,...,\pc} {
      \pgfmathtruncatemacro{\ex}{\r<=\d}
      \ifnum\ex=1 \def\cor{extra}\else\def\cor{white}\fi
      \draw[fill=\cor] ({\xo+\pl},{-\r}) rectangle ++(1,1);
    }
    % casa de partida (ponto) e casa de chegada (x)
    \fill ({\xo+0.5},0.5) circle (0.13);
    \node at ({\xo+\pl+0.5},{-\pc+0.5}) {$\times$};
    % comprimento de cada perna, em passos
    \node[gray!60!black] at ({\xo+(\pl+1)/2},1.45) {\pl};
    \node[gray!60!black] at ({\xo+\pl+1.5},{-\pc/2}) {\pc};
    % rótulo do dígito
    \node[anchor=north] at ({\xo+(\pl+1)/2},-3.4) {dígito $\d$};
  }
\end{tikzpicture}

  \caption{Forma do pulo para os dígitos 0, 1 e 2. As casas brancas formam o L do dígito 0, e as
  casas coloridas são as $d$ casas acrescentadas em cada perna. O ponto é a casa de partida, o
  $\times$ é a casa de chegada, e os números indicam o comprimento de cada perna, em passos.}
  \label{fig:pulos}
\end{figure}

O enunciado traz um tabuleiro de exemplo, com 10 linhas e 10 colunas
(Figura~\ref{fig:exemplo}), e informa que nele o cavalo vai de C a S em 3 pulos. Esse
exemplo é usado ao longo do relatório para explicar e conferir o algoritmo.

\begin{figure}[ht]
  \centering
  \ttfamily
  \begin{tabular}{@{}r@{\hspace{1.2em}}l@{}}
     & 0123456789\\
    0 & 4286305993\\
    1 & 67126514C1\\
    2 & 3541559238\\
    3 & 6344361236\\
    4 & 0419867802\\
    5 & 4271613350\\
    6 & S668962685\\
    7 & 6244051139\\
    8 & 3714702030\\
    9 & 0084825587\\
  \end{tabular}
  \caption{Tabuleiro de exemplo do enunciado. Linhas e colunas são numeradas de 0 a 9, a partir do
  canto superior esquerdo: C está na linha 1, coluna 8, e S na linha 6, coluna 0.}
  \label{fig:exemplo}
\end{figure}

Cada caso de teste é um arquivo de texto com $N$ linhas de $N$ caracteres: dígitos de 0 a 9,
exatamente um C e exatamente um S (isso foi conferido nos oito arquivos). Para cada caso, a saída
é o menor número de movimentos de C até S, ou a palavra ``impossível'' quando S não pode ser
alcançada. Os oito tabuleiros vão de $40 \times 40$ a $1500 \times 1500$ casas, e o maior
deles tem mais de dois milhões de casas.

O enunciado não exige recursão, e a solução apresentada usa um laço com uma fila, pelos
motivos explicados na Seção~\ref{sec:algoritmo}.

Este relatório está organizado da seguinte forma: a Seção~\ref{sec:modelagem} mostra como o
problema foi modelado, incluindo a regra do pulo e o tratamento do toro; a
Seção~\ref{sec:algoritmo} apresenta o algoritmo, com um exemplo passo a passo; a
Seção~\ref{sec:eficiencia} analisa a eficiência; a Seção~\ref{sec:resultados} mostra os
resultados dos oito casos e como eles foram conferidos; e a Seção~\ref{sec:conclusao} traz as
conclusões.

\section{Modelagem}
\label{sec:modelagem}

\subsection{O tabuleiro como um grafo}
\label{sec:grafo}

Para tratar o problema com algoritmos conhecidos, o tabuleiro é visto como um grafo. Cada
casa é um vértice, e há uma aresta da casa A para a casa B quando o cavalo consegue ir de A
até B em um pulo. Como cada pulo conta como um movimento, todas as arestas têm o mesmo custo,
1. Assim, o menor número de movimentos de C até S é o menor número de arestas de um caminho
de C até S. Esse é o problema do \emph{menor caminho em um grafo sem pesos}.

O grafo é dirigido: como o tamanho do pulo depende do dígito de uma casa, pode haver aresta de
A para B sem haver de B para A. Também não se constrói o grafo na memória. Ele é \emph{implícito}:
dada uma casa, calculam-se na hora as casas para onde ela leva, a partir do seu dígito. Isso
evita guardar as até $8N^2$ arestas, que são 18 milhões no maior tabuleiro.

A busca em largura (BFS, do inglês \emph{breadth-first search}) resolve exatamente esse
problema~\cite{cormen}. Ela visita as casas por camadas: primeiro a casa C, a distância 0;
depois todas as casas a 1 pulo de C; depois as casas a 2 pulos; e assim por diante. Uma fila
garante essa ordem, porque as casas descobertas a partir da camada $k$ entram na fila depois
de todas as casas da camada $k$, e portanto só são visitadas depois delas. Por isso, a
primeira vez que a busca chega a S é pelo caminho mais curto: se houvesse um caminho com
menos pulos, S estaria em uma camada anterior e já teria sido encontrada.

\subsection{O tabuleiro toroidal}
\label{sec:toro}

Para dar a volta nas bordas, usa-se aritmética modular. Se o cavalo está na linha $l$ e na
coluna $c$ e faz um pulo que soma $a$ às linhas e $b$ às colunas, ele chega à linha
$(l + a) \bmod N$ e à coluna $(c + b) \bmod N$. Isso vale mesmo quando o pulo é maior que o
tabuleiro inteiro.

Por exemplo, no tabuleiro $10 \times 10$ do enunciado, uma casa com dígito 9 faz um pulo de
pernas 10 e 11. Como $10 \bmod 10 = 0$ e $11 \bmod 10 = 1$, o cavalo volta para a mesma linha
e chega à coluna seguinte: dá uma volta completa no tabuleiro e termina ao lado de onde saiu.

\subsection{A regra do pulo}
\label{sec:regra}

O enunciado não dá uma fórmula para o pulo, só a figura com os L dos dígitos 0, 1 e 2
(Figura~\ref{fig:pulos}). Por isso é preciso extrair a regra da própria figura. Nela, o L do
dígito 0 tem pernas de 2 e 1 casas, o do dígito 1 tem pernas de 3 e 2, e o do dígito 2 tem
pernas de 4 e 3. As casas coloridas são o que foi acrescentado: $d$ casas em cada perna.
Conclui-se que um dígito $d$ forma um L de pernas $1 + d$ e $2 + d$. Combinando os sinais e
trocando qual perna vai na linha e qual vai na coluna, cada casa tem oito pulos possíveis:
$(\pm(1 + d), \pm(2 + d))$ e $(\pm(2 + d), \pm(1 + d))$. Com $d = 0$ recupera-se o pulo
normal de xadrez, como o enunciado pede.

Essa leitura foi conferida no exemplo do enunciado, em que C chega a S em 3 pulos. Como o dígito
da casa C está escondido pela letra, o exemplo foi executado com cada dígito possível, de 0 a 9,
no lugar do C. Com pernas $(1 + d, 2 + d)$ o resultado é sempre 3 pulos. Também foram testadas
duas outras regras, para ver se o exemplo e a figura as distinguem (Tabela~\ref{tab:hipoteses}).

\begin{table}[ht]
  \centering
  \caption{Regras testadas para as pernas do L, com o dígito escondido sob o C variando de 0 a 9.}
  \label{tab:hipoteses}
  \small
  \begin{tabular}{|c|c|c|}
    \hline
    \textbf{Pernas para o dígito $d$} & \textbf{Pulos de C a S no exemplo} & \textbf{Compatível com a figura?}\\
    \hline
    $(1 + d,\ 2 + d)$ & 3, com todos os dígitos & sim\\
    \hline
    $(2,\ 1 + d)$ ou $(1 + d,\ 2)$ & 3 com nove dígitos, 1 com o dígito 4 & não\\
    \hline
    $(1,\ 2 + d)$ ou $(2 + d,\ 1)$ & 2, 3 ou 4, conforme o dígito & não\\
    \hline
  \end{tabular}
\end{table}

A segunda regra não é descartada pelo exemplo, mas é descartada pela figura: para $d = 1$
ela daria pernas iguais, 2 e 2, que não formam um L. Na terceira, só uma perna cresce, o que
também contraria a figura, e o resultado no exemplo dependeria do dígito escondido. Só a
primeira regra é compatível com a figura e com o exemplo para qualquer dígito.

\subsection{Duas dúvidas de interpretação}
\label{sec:duvidas}

O texto do enunciado e o exemplo não resolvem dois pontos, e para cada um é preciso escolher
uma leitura.

\textbf{De qual casa vem o dígito.} O enunciado diz que o número de cada casa afeta o quanto
o cavalo pode pular. Adota-se a leitura mais direta: o cavalo usa o dígito da casa em que está
(leitura de \emph{origem}). A leitura alternativa, em que vale o dígito da casa onde ele vai
cair (leitura de \emph{destino}), também é compatível com o enunciado e também dá 3 pulos no
exemplo. Portanto o exemplo não serve para escolher entre as duas.

\textbf{O dígito escondido.} Nos arquivos, as letras C e S ocupam o lugar de um dígito, que se
perdeu. Na leitura de origem, isso afeta o primeiro pulo, que sai da casa C; a casa S não
importa, porque o cavalo para ao chegar nela. Como o enunciado define o dígito 0 como o pulo
normal de xadrez, supõe-se o dígito 0 sob o C, isto é, o primeiro pulo é um pulo normal de
cavalo. Na leitura de destino, o dígito escondido que importa é o de S, e faz-se a mesma
suposição.

\textbf{Quanto isso importa.} A Tabela~\ref{tab:sensibilidade} mostra o resultado de cada caso
nas quatro combinações: leitura de origem ou de destino, com o dígito escondido valendo 0 ou
sendo qualquer dígito de 0 a 9 (nesse caso, toma-se o menor resultado entre os dez). Só o
caso 1500 dá o mesmo resultado nas quatro combinações; nos outros sete, a resposta depende da
escolha.

\begin{table}[ht]
  \centering
  \caption{Menor número de movimentos de cada caso, segundo a leitura do enunciado e o dígito
  escondido. A primeira coluna é a leitura adotada neste relatório.}
  \label{tab:sensibilidade}
  \small
  \begin{tabular}{|c|c|c|c|c|}
    \hline
    & \multicolumn{2}{c|}{\textbf{Leitura de origem}} & \multicolumn{2}{c|}{\textbf{Leitura de destino}}\\
    \cline{2-5}
    \textbf{Caso} & \textbf{C = 0 (adotada)} & \textbf{C qualquer} & \textbf{S = 0} & \textbf{S qualquer}\\
    \hline
    40 & 4 & 4 & 2 & 2\\
    \hline
    80 & 6 & 4 & 6 & 4\\
    \hline
    100 & 8 & 6 & 8 & 6\\
    \hline
    150 & 11 & 9 & 11 & 11\\
    \hline
    200 & 9 & 9 & 11 & 9\\
    \hline
    400 & 19 & 17 & 17 & 17\\
    \hline
    800 & 43 & 43 & 43 & 41\\
    \hline
    1500 & 34 & 34 & 34 & 34\\
    \hline
  \end{tabular}
\end{table}

Por isso os números deste relatório são apresentados como os da leitura de origem com o dígito 0,
e a dependência fica registrada. Se o enunciado for esclarecido, o algoritmo da
Seção~\ref{sec:algoritmo} não muda: muda apenas a regra que diz para onde cada casa leva.
\pendente{Confirmar com o professor qual leitura vale. Se mudar, ajustar esta seção, o Resumo
e a Tabela~\ref{tab:resultados}.}

\section{Algoritmo}
\label{sec:algoritmo}

\subsection{Escolha do método}
\label{sec:metodo}

Como visto na Seção~\ref{sec:grafo}, o problema é achar um menor caminho em um grafo em que
todas as arestas custam 1. Três métodos clássicos foram considerados:

\begin{itemize}
  \item \textbf{Busca em profundidade (DFS).} Segue um caminho até o fim antes de tentar
    outro. Ela acha \emph{um} caminho até S, mas não garante que seja o mais curto; para
    garantir, seria preciso testar todos os caminhos possíveis, o que é inviável. Além disso, se
    fosse escrita com recursão, a profundidade poderia chegar a $N^2$ chamadas aninhadas (mais
    de dois milhões no maior tabuleiro), e a pilha de execução não suporta isso: no
    Trabalho~1, a pilha já estourou com algumas dezenas de milhares de chamadas.
  \item \textbf{Algoritmo de Dijkstra.} Acha menores caminhos quando as arestas têm custos
    diferentes, usando uma fila de prioridade. Aqui todos os custos são iguais, então esse
    trabalho extra é desnecessário: a BFS obtém o mesmo resultado com menos trabalho,
    $O(V + E)$ contra $O((V + E) \log V)$~\cite{cormen}. O Dijkstra é usado apenas como
    conferência independente dos resultados (Seção~\ref{sec:resultados}).
  \item \textbf{Busca em largura (BFS).} Explora por camadas, e uma fila simples basta. É o
    método adotado.
\end{itemize}

Como o enunciado não exige recursão, a BFS é escrita com um laço e uma fila, e não com
chamadas recursivas. Isso evita o limite da pilha e deixa o código mais direto: a própria fila
guarda o que ainda falta visitar.

\subsection{O algoritmo}
\label{sec:pseudocodigo}

O Algoritmo~\ref{alg:bfs} mostra a busca, e o Algoritmo~\ref{alg:pulos} mostra o cálculo dos
pulos de uma casa. Nas linhas, $(l, c)$ é uma casa (linha e coluna) e $N$ é o lado do
tabuleiro.

\begin{algorithm}[ht]
  \caption{Menor número de movimentos do cavalo de C até S}
  \label{alg:bfs}
  \begin{algorithmic}[1]
    \Function{MenorNumeroDeMovimentos}{$T$}
      \State $N \gets$ lado do tabuleiro $T$ \label{lin:lado}
      \State $(l_C, c_C) \gets$ posição de C em $T$;\quad $(l_S, c_S) \gets$ posição de S em $T$ \label{lin:pos}
      \State $dist[l][c] \gets -1$ para todas as casas \label{lin:init}
      \State $dist[l_C][c_C] \gets 0$ \label{lin:dist0}
      \State $fila \gets$ fila vazia;\quad insere $(l_C, c_C)$ na $fila$ \label{lin:fila0}
      \While{$fila$ não está vazia} \label{lin:while}
        \State $(l, c) \gets$ retira o primeiro elemento da $fila$ \label{lin:retira}
        \If{$(l, c) = (l_S, c_S)$} \label{lin:teste}
          \State \Return $dist[l][c]$ \label{lin:resposta}
        \EndIf
        \State $d \gets$ dígito da casa $(l, c)$ \Comment{0 se for a casa C} \label{lin:digito}
        \ForAll{$(a, b)$ em \Call{Pulos}{$d$}} \label{lin:for}
          \State $l' \gets (l + a) \bmod N$;\quad $c' \gets (c + b) \bmod N$ \label{lin:mod}
          \If{$dist[l'][c'] = -1$} \label{lin:novo}
            \State $dist[l'][c'] \gets dist[l][c] + 1$ \label{lin:marca}
            \State insere $(l', c')$ na $fila$ \label{lin:insere}
          \EndIf
        \EndFor
      \EndWhile
      \State \Return $-1$ \Comment{S não pode ser alcançada} \label{lin:imposs}
    \EndFunction
  \end{algorithmic}
\end{algorithm}

\begin{algorithm}[ht]
  \caption{Pulos possíveis a partir de uma casa com dígito $d$}
  \label{alg:pulos}
  \begin{algorithmic}[1]
    \Function{Pulos}{$d$}
      \State $p \gets 1 + d$;\quad $q \gets 2 + d$ \label{lin:pq}
      \State \Return $\{(\pm p, \pm q),\ (\pm q, \pm p)\}$ \Comment{8 pares, combinando os sinais} \label{lin:oito}
    \EndFunction
  \end{algorithmic}
\end{algorithm}

\textbf{Como funciona.} As linhas \ref{lin:lado} a \ref{lin:fila0} preparam a busca. A matriz
$dist$ guarda, para cada casa, quantos pulos foram necessários para chegar a ela, e o valor
$-1$ significa ``ainda não descoberta'' (linha~\ref{lin:init}). A casa C começa com distância 0
e é a única na fila (linhas~\ref{lin:dist0} e~\ref{lin:fila0}).

O laço das linhas \ref{lin:while} a \ref{lin:insere} é o coração do algoritmo. A cada volta,
ele retira da fila a casa mais antiga (linha~\ref{lin:retira}). Se for S, a distância dela é a
resposta (linhas~\ref{lin:teste} e~\ref{lin:resposta}): pelo argumento da
Seção~\ref{sec:grafo}, a primeira vez que S sai da fila já é pelo caminho mais curto. Se não
for, o algoritmo pega o dígito da casa (linha~\ref{lin:digito}) e percorre os 8 pulos que esse
dígito permite (linha~\ref{lin:for}). Para cada pulo, calcula a casa de chegada com aritmética
modular, que dá a volta no toro (linha~\ref{lin:mod}); se essa casa ainda não foi descoberta
(linha~\ref{lin:novo}), marca a distância dela, que é a da casa atual mais um
(linha~\ref{lin:marca}), e a coloca no fim da fila (linha~\ref{lin:insere}). Se a fila esvazia
sem que S apareça, todas as casas alcançáveis já foram visitadas e S não está entre elas:
a resposta é $-1$, que o programa apresenta como ``impossível'' (linha~\ref{lin:imposs}).

\subsection{Decisões de implementação}
\label{sec:decisoes}

\begin{itemize}
  \item \textbf{A matriz $dist$ faz dois papéis.} Ela guarda a distância de cada casa e, ao
    mesmo tempo, marca quais casas já foram descobertas ($-1$ quando não foram). Assim não
    é preciso uma segunda estrutura de visitadas, e o teste da linha~\ref{lin:novo} leva
    tempo constante.
  \item \textbf{A casa é marcada ao entrar na fila, não ao sair.} Se a marcação fosse só ao sair,
    uma mesma casa poderia entrar na fila várias vezes, uma para cada vizinho que a descobrisse,
    e o trabalho aumentaria. Marcando na linha~\ref{lin:marca}, cada casa entra no máximo uma
    vez, e é isso que limita o custo total (Seção~\ref{sec:eficiencia}).
  \item \textbf{A fila é um \texttt{ArrayDeque} do Java}, que insere no fim e retira do início
    em tempo constante.
  \item \textbf{Usa-se \texttt{Math.floorMod} em vez do operador \texttt{\%}.} Em Java,
    \texttt{\%} devolve resultado negativo para números negativos (por exemplo,
    $-2 \bmod 10$ dá $-2$, e não 8), e as coordenadas de chegada podem ficar negativas quando
    o pulo vai para cima ou para a esquerda. O \texttt{floorMod} devolve sempre um valor entre
    0 e $N - 1$.
  \item \textbf{O dígito da casa C (linha~\ref{lin:digito}) vale 0}, conforme a regra do pulo
    (Seção~\ref{sec:regra}). A casa S nunca é expandida, porque o laço termina ao retirá-la.
  \item \textbf{A busca termina ao retirar S da fila (linha~\ref{lin:teste}).} Poderia terminar
    antes, ao descobrir S; esse ganho foi medido na Seção~\ref{sec:conclusao}.
\end{itemize}

\subsection{Exemplo passo a passo}
\label{sec:exemplo}

O algoritmo foi aplicado ao tabuleiro da Figura~\ref{fig:exemplo}. A BFS descobre as casas em
camadas: 1 casa a 0 pulos (a própria C), 8 casas a 1 pulo, 36 a 2 pulos, 42 a 3 pulos e 13 a
4 pulos, num total de 100, ou seja, todas as casas do tabuleiro são alcançáveis. A casa S
está na camada 3. Ela é retirada da fila depois que 67 casas foram retiradas (contando a
própria S), e a busca termina.

Para saber o caminho, anda-se de S para trás: em cada camada procura-se uma casa cujo pulo
chega na casa seguinte. Um dos caminhos mínimos é
$C(1, 8) \to (2, 6) \to (2, 7) \to S(6, 0)$, detalhado na Tabela~\ref{tab:trace}.

\begin{table}[ht]
  \centering
  \caption{Um caminho mínimo no tabuleiro de exemplo, pulo a pulo. As coordenadas são
  (linha, coluna) e o tabuleiro tem $N = 10$.}
  \label{tab:trace}
  \small
  \begin{tabular}{|c|c|c|c|c|c|c|}
    \hline
    \textbf{Pulo} & \textbf{Sai de} & \textbf{Dígito} & \textbf{Pernas} & \textbf{Deslocamento}
      & \textbf{Sem o módulo} & \textbf{Chega em}\\
    \hline
    1 & $(1, 8)$ & 0 & $(1, 2)$ & $(+1, -2)$ & $(2, 6)$ & $(2, 6)$\\
    \hline
    2 & $(2, 6)$ & 9 & $(10, 11)$ & $(+10, +11)$ & $(12, 17)$ & $(2, 7)$\\
    \hline
    3 & $(2, 7)$ & 2 & $(3, 4)$ & $(+4, +3)$ & $(6, 10)$ & $(6, 0)$\\
    \hline
  \end{tabular}
\end{table}

O primeiro pulo sai de C e, pela regra da Seção~\ref{sec:regra}, é um pulo normal de
cavalo. O segundo mostra o toro em ação: o dígito 9 dá pernas de 10 e 11, que em um tabuleiro
de lado 10 equivalem a 0 e 1, e o cavalo chega à casa vizinha da direita. O terceiro sai pela
borda direita (coluna 10) e entra pela esquerda (coluna 0), caindo em S. A
Figura~\ref{fig:camadas} mostra as camadas e o caminho no tabuleiro.

\begin{figure}[ht]
  \centering
  % Figura 3: camadas da BFS no tabuleiro de exemplo e um caminho mínimo de C a S.
%
% Cada casa é dada por um par "dígito/distância": o dígito é o do tabuleiro (x marca as
% casas C e S, desenhadas à parte) e a distância é a do mapa impresso por
% MedicoesHiperpulos (ver Trabalho2/medicoes-referencia.txt). A cor da casa vem da distância.
% Para mudar as cores, edite os \colorlet; para mover a legenda, mude o valor de \xl.
\begin{tikzpicture}[x=0.82cm,y=0.82cm,font=\footnotesize]
  \colorlet{nivel0}{black}
  \colorlet{nivel1}{blue!10}
  \colorlet{nivel2}{blue!25}
  \colorlet{nivel3}{blue!40}
  \colorlet{nivel4}{blue!58}
  \colorlet{trilha}{orange!85!black}

  % tabuleiro: uma lista por linha, de cima para baixo
  \foreach \lin [count=\r from 0] in {
    {4/1,2/2,8/3,6/2,3/3,0/4,5/1,9/2,9/3,3/4},
    {6/2,7/3,1/2,2/3,6/4,5/3,1/2,4/3,x/0,1/3},
    {3/1,5/2,4/3,1/2,5/3,5/2,9/1,2/2,3/3,8/4},
    {6/2,3/3,4/2,4/3,3/2,6/3,1/2,2/1,3/4,6/1},
    {0/3,4/4,1/3,9/2,8/3,6/2,7/3,8/4,0/3,2/2},
    {4/2,2/3,7/2,1/3,6/2,1/3,3/2,3/3,5/4,0/3},
    {x/3,6/2,6/3,8/2,9/3,6/2,2/3,6/2,8/3,5/2},
    {6/2,2/3,4/2,4/3,0/2,5/3,1/2,1/3,3/2,9/3},
    {3/3,7/2,1/3,4/4,7/3,0/2,2/3,0/2,3/3,0/4},
    {0/2,0/3,8/4,4/3,8/2,2/3,5/4,5/1,8/4,7/1}%
  } {
    \foreach \dig/\dis [count=\c from 0] in \lin {
      \draw[fill=nivel\dis] (\c,{-\r-1}) rectangle ++(1,1);
      \ifx\dig x\else \node[black!75] at ({\c+0.5},{-\r-0.5}) {\dig}; \fi
    }
  }

  % casa C (partida) e casa S (saída)
  \node[white,font=\footnotesize\bfseries] at (8.5,-1.5) {C};
  \draw[trilha,very thick] (0,-7) rectangle ++(1,1);
  \node[font=\footnotesize\bfseries] at (0.5,-6.5) {S};

  % numeração de linhas e colunas
  \foreach \i in {0,...,9} {
    \node[gray!50!black] at ({\i+0.5},0.4) {\i};
    \node[gray!50!black] at (-0.4,{-\i-0.5}) {\i};
  }

  % caminho mínimo: C -> (2,6) -> (2,7) -> S
  \tikzset{seta/.style={trilha,very thick,-{Stealth[length=2.2mm]},shorten >=3pt,shorten <=3pt},
           num/.style={circle,fill=white,draw=trilha,inner sep=1pt,font=\scriptsize\bfseries,text=trilha}}
  \draw[trilha,very thick] (6,-3) rectangle ++(2,1);  % casas (2,6) e (2,7), para o digito ficar legivel
  \draw[seta] (8.5,-1.5) -- (6.5,-2.5) node[num,midway] {1};
  \draw[seta,shorten >=6pt] (6.5,-2.5) -- (7.5,-2.5) node[num,midway,above=3pt] {2};
  % o pulo 3 sai pela borda direita (coluna 10) e entra pela esquerda (coluna 0)
  \draw[trilha,very thick,shorten <=3pt] (7.5,-2.5) -- (10,-5.833) node[num,pos=0.45] {3};
  \draw[seta,shorten <=0pt,shorten >=9pt] (0,-5.833) -- (0.5,-6.5);
  \fill[trilha] (10,-5.833) circle (0.1);   % ponto onde o pulo sai do tabuleiro
  \fill[trilha] (0,-5.833) circle (0.1);    % ponto onde ele volta a entrar

  % legenda
  \def\xl{10.7}
  \node[anchor=west,align=left] at (\xl,-0.6) {\textbf{Distância até C}\\(em pulos)};
  \foreach \k in {0,...,4} {
    \draw[fill=nivel\k] (\xl,{-2.2-0.85*\k}) rectangle ++(0.7,0.7);
    \node[anchor=west] at ({\xl+0.9},{-1.85-0.85*\k}) {\k};
  }
  \node[anchor=west,align=left,black!75] at (\xl,-7.0) {Número na casa:\\dígito do tabuleiro};
  \node[num] at ({\xl+0.25},-8.3) {1};
  \node[anchor=west,align=left,black!75] at ({\xl+0.6},-8.3) {ordem dos pulos};
\end{tikzpicture}

  \caption{Camadas da BFS no tabuleiro de exemplo e um caminho mínimo de C a S. A cor indica a
  distância até C, e as setas numeradas são os três pulos da Tabela~\ref{tab:trace}. O pulo 3
  sai pela borda direita e entra pela esquerda.}
  \label{fig:camadas}
\end{figure}

\section{Eficiência}
\label{sec:eficiencia}

\subsection{Análise do custo}
\label{sec:custo}

Seja $N$ o lado do tabuleiro. Há $V = N^2$ casas, e cada casa tem no máximo 8 pulos, então o
grafo tem $E \le 8N^2$ arestas.

\textbf{Tempo.} O trabalho do Algoritmo~\ref{alg:bfs} é contado por partes:

\begin{itemize}
  \item A inicialização de $dist$ (linha~\ref{lin:init}) percorre todas as casas:
    $\Theta(N^2)$.
  \item Cada casa entra na fila no máximo uma vez, porque só entra se $dist$ ainda vale $-1$
    (linha~\ref{lin:novo}) e é marcada na hora (linha~\ref{lin:marca}). Logo, o laço da
    linha~\ref{lin:while} roda no máximo $N^2$ vezes.
  \item A cada volta, o laço da linha~\ref{lin:for} executa exatamente 8 vezes (o
    Algoritmo~\ref{alg:pulos} sempre devolve 8 pares), e cada execução faz um número constante
    de operações: duas somas, dois módulos, uma consulta e uma atualização em $dist$ e uma
    inserção na fila, todas em tempo constante.
\end{itemize}

Somando, o laço principal faz no máximo $N^2 \times 8 \times c$ operações, para uma constante
$c$, ou seja, $O(N^2)$. Em termos de vértices e arestas, é o custo $\Theta(V + E)$ conhecido da
BFS~\cite{cormen}. Como a inicialização já custa $N^2$ em qualquer tabuleiro, o tempo do
algoritmo é $\Theta(N^2)$ sempre. O que muda de um tabuleiro para outro é a constante: quando S
é encontrada cedo, a busca visita só parte das casas.

\textbf{Espaço.} A matriz $dist$ ocupa $N^2$ inteiros, a fila guarda no máximo $N^2$ casas
(cada uma entra uma vez) e o próprio tabuleiro ocupa $N^2$ caracteres. O espaço é, portanto,
$\Theta(N^2)$.

Não dá para melhorar a \emph{ordem de grandeza}: só para localizar S no tabuleiro, qualquer
algoritmo precisa, no pior caso, olhar as $N^2$ casas. A BFS é, então, assintoticamente ótima
para este problema, e os ganhos que restam são de constante (Seção~\ref{sec:conclusao}).

\subsection{Conferência com as medições}
\label{sec:medicoes}

Se o custo é proporcional a $N^2$, dobrar $N$ deve multiplicar o tempo por 4. Para conferir,
mede-se o pior caso: uma busca que não para em S e visita todas as casas alcançáveis. Isso
isola o efeito do tamanho do tabuleiro, sem depender de onde S está.

As medições foram feitas em um Apple M2 com 16~GB de memória, macOS 26.7 e Java 24.0.1, sem contar a
leitura do arquivo. Cada tempo é o melhor de 15 execuções na mesma JVM, o que descarta o custo de
compilação das primeiras execuções e a interferência da coleta de lixo. Também se fixou o
tamanho da memória da JVM (\texttt{-Xms2g -Xmx2g}): com tamanho automático, o tempo absoluto
variou em até 1,8 vez de uma execução para outra, mas as razões entre tamanhos, que são o que
se usa aqui, não mudaram.

\begin{table}[ht]
  \centering
  \caption{Exploração completa, sem parar em S. As duas últimas colunas comparam cada tamanho
  com o anterior: se o custo é proporcional a $N^2$, elas devem ser parecidas.}
  \label{tab:crescimento}
  \small
  \begin{tabular}{|c|r|r|c|c|c|}
    \hline
    \textbf{$N$} & \textbf{Casas alcançadas} & \textbf{Tempo (ms)} & \textbf{ns por casa}
      & \textbf{Tempo / anterior} & \textbf{$N^2$ / anterior}\\
    \hline
    40 & 1.600 & 0,14 & 88 & - & -\\
    \hline
    80 & 6.398 & 0,62 & 96 & 4,4 & 4,0\\
    \hline
    100 & 9.999 & 0,97 & 97 & 1,6 & 1,6\\
    \hline
    150 & 22.494 & 2,23 & 99 & 2,3 & 2,3\\
    \hline
    200 & 39.992 & 4,02 & 101 & 1,8 & 1,8\\
    \hline
    400 & 159.950 & 16,41 & 103 & 4,1 & 4,0\\
    \hline
    800 & 639.857 & 67,35 & 105 & 4,1 & 4,0\\
    \hline
    1500 & 2.249.516 & 241,98 & 108 & 3,6 & 3,5\\
    \hline
  \end{tabular}
\end{table}

A Tabela~\ref{tab:crescimento} confirma a previsão: a razão entre os tempos acompanha a razão
entre os valores de $N^2$ em todos os tamanhos. O tempo por casa visitada, que deveria ser
constante, sobe pouco, de 88 para 108 ns (23\%), enquanto o número de casas cresce 1400 vezes.
Uma hipótese para essa pequena alta é o uso de memória: o maior tabuleiro ocupa dezenas de
megabytes, que não cabem nos caches do processador; isso não foi medido. Para $N \le 100$ os
tempos ficam abaixo de 1~ms, perto do limite do que o cronômetro e a JVM medem com precisão, e
as razões desses tamanhos devem ser lidas com cuidado.

\begin{figure}[ht]
  \centering
  \begin{tikzpicture}
    \begin{loglogaxis}[
        width=0.9\linewidth, height=7.5cm,
        xlabel={Lado do tabuleiro $N$ (casas)},
        ylabel={Tempo (ms)},
        xmin=30, xmax=2000, ymin=0.01, ymax=500,
        xtick={40,100,200,400,800,1500}, xticklabels={40,100,200,400,800,1500},
        ytick={0.01,0.1,1,10,100}, yticklabels={0{,}01,0{,}1,1,10,100},
        xminorticks=false, yminorticks=false,
        grid=major, grid style={line width=0.2pt,draw=gray!40},
        tick label style={font=\small}, label style={font=\small},
        legend pos=north west,
        legend style={font=\small,draw=none,fill=white,fill opacity=0.85,text opacity=1},
      ]
      \addplot[mark=*,thick,blue] coordinates {
        (40,0.14) (80,0.62) (100,0.97) (150,2.23) (200,4.02) (400,16.41) (800,67.35) (1500,241.98)
      };
      \addlegendentry{Exploração completa (pior caso)}
      \addplot[only marks,mark=square*,red] coordinates {
        (40,0.02) (80,0.40) (100,0.78) (150,1.92) (200,1.50) (400,8.79) (800,57.24) (1500,36.30)
      };
      \addlegendentry{Solver real (para em S)}
      \addplot[dashed,black,domain=40:1500,samples=2] {241.98*(x/1500)^2};
      \addlegendentry{Referência proporcional a $N^2$}
    \end{loglogaxis}
  \end{tikzpicture}
  \caption{Tempo de execução por tamanho do tabuleiro, em escala logarítmica nos dois eixos. A
  exploração completa acompanha a reta tracejada, proporcional a $N^2$, que passa pelo ponto
  $N = 1500$. Os pontos vermelhos são os casos reais, em que a busca para ao chegar a S.}
  \label{fig:crescimento}
\end{figure}

\subsection{Por que os casos reais não crescem de forma regular}
\label{sec:casos-reais}

Nos casos reais o algoritmo para ao retirar S da fila, e o trabalho depende de onde S está em
relação a C. A Tabela~\ref{tab:resultados} mostra que a busca retira da fila desde 16,7\% das
casas (caso 1500) até 93,3\% (caso 800). Por isso o caso 1500, que tem 3,5 vezes mais casas
que o caso 800, é resolvido mais rápido (36,30~ms contra 57,24~ms): a BFS chega a S depois de
visitar uma parte pequena do tabuleiro, e o resto nem chega a ser explorado. É por isso que os
pontos vermelhos da Figura~\ref{fig:crescimento} ficam abaixo da curva do pior caso e não
seguem uma reta.

\section{Resultados}
\label{sec:resultados}

Aplicando o algoritmo da Seção~\ref{sec:algoritmo} aos oito casos de teste, com a regra do
pulo da Seção~\ref{sec:regra}, obtêm-se os resultados da Tabela~\ref{tab:resultados}.

\begin{table}[ht]
  \centering
  \caption{Resultados dos oito casos de teste. A coluna ``Casas retiradas'' conta as casas que
  saíram da fila até a retirada de S, e o tempo é o melhor de 15 execuções.}
  \label{tab:resultados}
  \small
  \begin{tabular}{|c|c|c|r|c|r|}
    \hline
    \textbf{Caso} & \textbf{$N$} & \textbf{Movimentos} & \textbf{Casas retiradas}
      & \textbf{\% de $N^2$} & \textbf{Tempo (ms)}\\
    \hline
    caso40 & 40 & 4 & 425 & 26,6 & 0,02\\
    \hline
    caso80 & 80 & 6 & 4.483 & 70,0 & 0,40\\
    \hline
    caso100 & 100 & 8 & 8.596 & 86,0 & 0,78\\
    \hline
    caso150 & 150 & 11 & 20.863 & 92,7 & 1,92\\
    \hline
    caso200 & 200 & 9 & 15.483 & 38,7 & 1,50\\
    \hline
    caso400 & 400 & 19 & 92.523 & 57,8 & 8,79\\
    \hline
    caso800 & 800 & 43 & 597.282 & 93,3 & 57,24\\
    \hline
    caso1500 & 1500 & 34 & 375.230 & 16,7 & 36,30\\
    \hline
  \end{tabular}
\end{table}

Nenhum dos oito casos é impossível: em todos, S é alcançável a partir de C. O número de
movimentos vai de 4 a 43, e o tabuleiro de maior lado, $1500 \times 1500$, foi resolvido em
36~ms. O caso mais demorado é o 800, com 57~ms, e a razão disso, que não é o tamanho, está na
Seção~\ref{sec:casos-reais}.

\subsection{Como os resultados foram conferidos}
\label{sec:conferencia}

Para não confiar em um único programa, os resultados foram conferidos de várias formas:

\begin{itemize}
  \item \textbf{Exemplo do enunciado.} O algoritmo dá 3 pulos, como o enunciado informa, e uma
    busca exaustiva confirma que 2 pulos não bastam.
  \item \textbf{Segunda implementação.} A busca foi escrita de novo, separadamente (casas
    guardadas como números inteiros e fila em vetor), para que um erro do primeiro programa não
    se repetisse igual na conferência. As duas dão a mesma resposta nos oito casos.
  \item \textbf{Algoritmo de Dijkstra.} Um método diferente, com fila de prioridade, dá a
    mesma distância nos casos até $400 \times 400$.
  \item \textbf{Caminho legal.} Reconstruiu-se o caminho mínimo de cada caso e conferiu-se, pulo
    a pulo, que todos obedecem à regra do pulo.
  \item \textbf{Tabuleiros pequenos aleatórios.} Em 389 tabuleiros sorteados (com semente
    fixa), a distância encontrada é atingível e uma busca exaustiva em profundidade limitada
    não acha caminho menor.
  \item \textbf{Casos de borda.} Tabuleiros construídos à mão testam o pulo padrão, um pulo
    que só fecha dando a volta no toro e o dígito 9 em um tabuleiro $7 \times 7$ (pernas
    maiores que o lado, o que testa o módulo). A saída ``impossível'' foi exercitada com um
    tabuleiro sem solução, e a busca independente também não achou caminho.
\end{itemize}

Essas conferências mostram que o programa implementa corretamente a regra do pulo da
Seção~\ref{sec:regra}. Como as implementações independentes seguem a mesma regra, elas testam a
implementação, e não a interpretação do enunciado.

\section{Conclusão}
\label{sec:conclusao}

O problema foi resolvido modelando-se o tabuleiro como um grafo implícito, com uma casa por vértice e
um pulo por aresta, e usando busca em largura. A modelagem transformou ``menor número de
movimentos'' em ``menor caminho em um grafo sem pesos'', problema para o qual a BFS é correta e
custa $\Theta(N^2)$, o que é ótimo em ordem de grandeza. Os oito casos foram resolvidos, com 4
a 43 movimentos e no máximo 57~ms, e os resultados foram conferidos de várias formas
(Seção~\ref{sec:conferencia}). As medições confirmaram o custo quadrático: dobrar $N$ multiplica
o tempo por cerca de 4 (Tabela~\ref{tab:crescimento}).

\subsection{Melhorias possíveis}
\label{sec:melhorias}

Mesmo sendo ótimo em ordem de grandeza, o algoritmo pode ficar mais rápido por um fator
constante. Duas melhorias foram implementadas e medidas, sempre conferindo que a resposta continua a
mesma nos oito casos e no exemplo.

\textbf{Parar ao descobrir S (versão B).} A BFS atribui a uma casa a sua distância definitiva
no momento em que a descobre. Portanto não é preciso esperar S chegar ao início da fila para
devolver a resposta, como é feito na linha~\ref{lin:teste}: basta testar ao inserir
(linha~\ref{lin:insere}). Isso evita expandir as casas da última camada que ainda estão na fila
antes de S. A versão B foi de 1,5 a 1,9 vez mais rápida que a original nos casos com $N \ge 100$
(Tabela~\ref{tab:melhorias}).

\textbf{Casas como inteiros e pulos sem alocação (versão C).} Na versão original, cada casa
visitada aloca cerca de dez vetores pequenos: um par $(l, c)$ para entrar na fila e os 8 pares
de deslocamento devolvidos por \textsc{Pulos}. A versão C, que também para ao descobrir S,
guarda a casa como um único inteiro $l \times N + c$, usa um vetor de inteiros como fila e
calcula os 8 pulos no próprio laço, sem criar objetos. Ela foi de 2,9 a 4,4 vezes mais rápida que
a original. Acredita-se que o ganho venha de evitar as alocações e de acessar a memória de forma
mais compacta, mas cada efeito não foi isolado.

\begin{table}[ht]
  \centering
  \caption{Tempo das três versões, em ms (melhor de 15 execuções), e o quanto a versão original
  (A) é mais lenta que as outras. Os dois menores casos ($N = 40$ e $N = 80$) ficam abaixo de 0,5~ms e
  foram omitidos pelo ruído de medição. A medição é independente da Tabela~\ref{tab:resultados}, então a
  versão A difere dela em alguns por cento.}
  \label{tab:melhorias}
  \small
  \begin{tabular}{|c|c|c|c|c|c|}
    \hline
    \textbf{$N$} & \textbf{A: original} & \textbf{B: para ao descobrir S}
      & \textbf{C: inteiros e B} & \textbf{A / B} & \textbf{A / C}\\
    \hline
    100 & 0,75 & 0,39 & 0,17 & 1,9 & 4,4\\
    \hline
    150 & 1,94 & 1,09 & 0,48 & 1,8 & 4,0\\
    \hline
    200 & 1,36 & 0,75 & 0,34 & 1,8 & 4,0\\
    \hline
    400 & 8,77 & 5,12 & 2,43 & 1,7 & 3,6\\
    \hline
    800 & 57,95 & 38,26 & 16,39 & 1,5 & 3,5\\
    \hline
    1500 & 36,36 & 23,25 & 12,61 & 1,6 & 2,9\\
    \hline
  \end{tabular}
\end{table}

Duas outras ideias não foram implementadas. A primeira é a busca bidirecional, a partir de C e
de S ao mesmo tempo. Aqui ela não é imediata: como o pulo depende da casa de onde se sai,
andar para trás a partir de S exige saber quais casas chegam a cada casa, e isso obriga a montar
o grafo inverso, que custa $\Theta(N^2)$ de tempo e até $8N^2$ entradas de memória. Não se sabe
se o ganho compensaria esse custo; seria preciso medir. A segunda vale se o mesmo tabuleiro
tivesse várias saídas: a exploração completa da Seção~\ref{sec:medicoes} já calcula a distância
de C a todas as casas, então uma única busca responderia a todas.

Por fim, o que mais pesou neste trabalho não foi o algoritmo, que é padrão, e sim a modelagem:
decidir o que é vértice e o que é aresta, como tratar o toro e como extrair da figura do
enunciado a regra do pulo. Foram essas decisões que definiram os resultados, e por isso buscou-se
registrar cada uma com os seus motivos.


\begin{thebibliography}{1}
\bibitem{cormen}
Cormen, T. H.; Leiserson, C. E.; Rivest, R. L.; Stein, C.:
\textbf{Introduction to Algorithms}. 3.\,ed. MIT Press, Cambridge, 2009.
(Busca em largura: seção 22.2; algoritmo de Dijkstra: seção 24.3.)
\end{thebibliography}

\end{document}
